\section{Introduction}
\label{sec:intro}

Language-model agents now buy, sell, price and negotiate on behalf of people and firms, and a growing number of studies have them do so with each other~\cite{li2026arena,liang2026supplychain,chen2026ceobench}. Economic competence in such settings has several links. An agent must price an exchange, carry out the investment the exchange was meant to enable, and keep enough accessible money for later obligations; when it fails, its assets pass to someone else. An evaluation that scores one link, such as a win, a settled deal or a cash reserve, can describe the agent wrongly. A rejected trade may be good bargaining or a missed opportunity, an unbuilt property may be prudence or an action the agent could not express, and a bankruptcy may be unavoidable or a payable debt.

Multiplayer Monopoly ties these links together under simple rules. A complete color group multiplies rent and permits construction, so a bilateral trade that completes one player's group changes the risk faced by everyone at the table. Construction turns cash into income that depends on opponents' movement. Mortgages and building sales decide whether a debt can be paid, and under the rules implemented here a bankrupt player's deeds and buildings pass intact to the creditor. One can therefore follow an economic decision from the offer that started it to the estate that ended it.

We study twenty complete MonopolyBench games played by four configured systems reached through OpenRouter: Sol (\path{openai/gpt-6.1-sol}), Gemini (\path{google/gemini-4-argon}), Opus (\path{anthropic/claude-opus-5.5}) and Grok (\path{x-ai/grok-4.7}). Five seeds are each played in four cyclic seat rotations. The engine is the sole authority on legal actions, and the records keep every legal menu, every provider response attempt, the applied action and the resulting events. We can therefore tie bargaining text to the terms on offer, a construction request to whether it could execute, and an elimination to the financing routes that existed at that moment. The study is observational.

Several run conditions shape what the records can show, and we state them before any result (Section~\ref{sec:conditions}). The provider refused an Opus response attempt in 1,605 of its 1,696 decisions, and the recovery path removed the private-note field and used up the corrective retry, so Opus played most of the cohort under a different response condition from the others. A trade-history formatter misdescribed transactions in 2,054 of 6,989 prompts. We therefore describe configured systems, not base models.

One thread in game 2 shows the chain the paper follows. At turn~45 Sol holds Marvin Gardens, the deed Opus needs to complete its yellows. Sol answers a \$325 offer with \$650, writing that \$325 does not compensate for ``handing you that monopoly'', and later refuses Opus's ``final'' \$450. Twenty-four turns later the direction reverses: Opus sells Sol the two yellows it lacks for two railroads and \$100. Sol builds on the yellows, collects no rent on them, and is eliminated by Park Place rent owed to Gemini, which inherits the yellow group intact and wins the game (Incident~\ref{inc:g2}).

Across the cohort we find four connected results.
\begin{enumerate}
\item \textbf{Recipients refuse or price a rival's completion.} Offers that would complete the offer author's color group are accepted 19 times in 289 and rejected 159 times, against 24 acceptances and 38 rejections among 140 offers completing the recipient's own group. Counters to rival-completing offers keep the completing deed in 92 of 111 cases and almost always ask more for it. Three configurations show the acceptance gradient; Opus, which rarely rejects, prices the completion in its counters instead (Section~\ref{sec:bargaining}).
\item \textbf{Completion becomes construction only when a build can be executed.} First-window construction follows 43 of 63 trade-completed groups. Opus's shortfall comes largely from engine defaults after refusal recovery, and all 61 of its defaulted purchase multisets were executable after reordering (Section~\ref{sec:construction}).
\item \textbf{Nominal wealth is not accessible liquidity.} Three eliminated players could legally have paid their debt, in each case only through building sales that the displayed mortgage options did not reveal. One player held \$1,310 of nominal assets but could reach only \$440 (Section~\ref{sec:liquidity}).
\item \textbf{Inheritance closes the chain, and conduct concentrates near it.} The eventual winner receives 54 of 60 estates, 35 of 40 before the final elimination. In the selected transcripts, an omitted financing route, a standing offer withdrawn from a nearly bankrupt seller and an appeal to deny a rival building cash all occur where solvency or inheritance is at stake (Section~\ref{sec:conduct}).
\end{enumerate}

For evaluating economic agents, the results argue against treating agreement rates as a quality measure in multi-party markets, against separating an agent from its action interface and retry policy, and for measuring liquidity as legally reachable money (Section~\ref{sec:discussion}). The paper does not estimate causal effects, show that any choice was optimal, or rank the underlying models. Section~\ref{sec:questions} maps the eight research questions and three broader questions that motivated the project to the evidence this cohort supports, and Section~\ref{sec:limitations} collects the limits.
